\subsection{几乎复流形与复流形}

本节假设 K = R，并记 X 为一个（实）\(\mathbf{C}^{r}\) 类流形，其中 \(r \in \mathbf{N}_{\mathbf{R}}\)。

8.8.1 (“复向量丛”). 以 X 为基底的 \(\mathbf{C}\) 上的向量丛（7.3.4）也称为以 X 为基底的复向量丛。若 M 是这样的向量丛，则称 M 的一个复向量坐标图为三元组 (U, \(\varphi\) , E)，其中 U 是 X 的开子集，E 是复 Banach 空间，且 \(\varphi\) 是从 \(M|_U\) 到平凡丛 \(E_U = U \times E\) 的复向量丛同构。其定义域 \((U_i,\varphi_i,E_i)\) 覆盖 X 的 M 的复向量坐标图族 \(U_i\) 称为 M 的复向量图册；每个复向量丛都具有一个复向量图册（7.3.3）。

有时用 C-向量代替复向量。同样，在谈论 R 上的向量丛、该丛的普通向量坐标图等对象时，有时用“实向量”或“R-向量”代替“向量”。

设 M 是以 X 为基底的复向量丛。M 的（实）底层向量丛（7.3.3）存在且仅存在一个自同构 j，其在每个纤维上的限制是乘以 C 的元素 i；j 的平方等于 -1。反之，若 M 是以 X 为基底的（实）向量丛，且 j 是 M 的满足 \(j^{2} = -1\) 的自同构，则 M 上存在唯一的复向量丛结构，以 M 为其（实）底层向量丛，并且其中 j 是乘以 i 的映射。

8.8.2（“向量丛的复化”）。对每个（实）Banach 空间 \(\tau\)，令 \(\tau(E) = E \otimes C\)，并对（实）Banach 空间的每个态射 \(E\)，令 \(\tau(u) = u \otimes Id_{C}\)，由此定义一个向量函子 \(u: E \to E'\)。

若 F 是以 X 为基底的 \(C^{k}\) 类向量丛（\(0 \leqslant k \leqslant r\)），则其在向量函子 \(\tau\) 下的像记为 \(F \otimes C\)。在 \(F \otimes C\) 上存在唯一的、以 X 为基底且属于 \(C^{k}\) 类的复向量丛结构；它与其（实）向量丛结构相容，并在每个纤维 \((\mathbf{F} \otimes \mathbf{C})_{x} = \mathbf{F}_{x} \otimes \mathbf{C}\) 上赋予其典范的复 Banach 空间结构（EVT, II, § 8, n° 1, Example）。赋予此结构后，称 \(F \otimes C\) 为 F 的复化。

设 U 是 X 的开子集；典范映射

\[
\mathcal {S} _ {\mathrm{F}} ^ {k} (\mathrm{U}) \otimes \mathbf {C} \rightarrow \mathcal {S} _ {\mathrm{F} \otimes \mathbf {C}} ^ {k} (\mathrm{U})
\]

是一个同构，借此将这两个空间等同。若 \(s \in \mathcal{S}_{\mathrm{F} \otimes \mathrm{C}}^{k}(\mathrm{U})\)，\(\sigma, \tau\) 是 \(\mathcal{S}_{\mathrm{F}}^{k}(\mathrm{U})\) 的元素且满足 \(s = \sigma + i\tau\)，则称它们为 \(s\) 的实部和虚部；截面 \(\bar{s} = \sigma - i\tau\) 称为 \(s\) 的共轭。

特别地，丛 \(\mathrm{T}(\mathrm{X}) \otimes \mathbf{C}\) 的截面称为 X 上的复向量场。

以 X 为底的复向量丛 H 的典范同构 \(\operatorname{Alt}_{\mathbf{R}}^{p}(\mathrm{T}_{x}(\mathrm{X});\mathrm{H})\to\operatorname{Alt}_{\mathbf{C}}^{p}(\mathrm{T}_{x}(\mathrm{X})\otimes\mathbf{C};\mathbf{H})\) 定义了从 \(\operatorname{Alt}_{\mathbf{R}}^{p}(\mathrm{T}(\mathrm{X});\mathrm{H})\) 到 \(\operatorname{Alt}_{\mathbf{C}}^{p}(\mathrm{T}(\mathrm{X})\otimes\mathbf{C};\mathbf{H})\) 的复向量丛典范同构，称为典范同构（7.8.5）；借此将这两个丛等同。设 \(\omega\) 是该丛的一个截面，换言之，是 X 上取值于 H 的 p 次微分形式；设 \(\zeta\) 是复向量场，其实部为 \(\xi\)，虚部为 \(\eta\)。置：

\[
i (\zeta , \omega) = i (\xi , \omega) + i. i (\eta , \omega).
\]

若 H 是平凡丛，置

\[
\theta_ {\zeta}. \omega = \theta_ {\xi}. \omega + i \theta_ {\eta}. \omega .
\]

同样，按线性将括号延拓到复向量场。前面各号的公式仍然成立。

8.8.3. X 上的一个 \(C^{k}\) 类几乎复结构（\(k \leqslant r - 1\)）是指在 \(\mathrm{T}(\mathrm{X})\) 上给定一个 \(C^{k}\) 类复向量丛结构，其底层实结构是 \(\mathrm{T}(\mathrm{X})\) 的自然结构。这种结构等价于给定 \(j\) 的一个 \(C^{k}\) 类自同构 \(\mathrm{T}(\mathrm{X})\)，满足 \(j^{2} = -1\)。

取这样一个结构。将 \(j\) 延拓为 \(\mathbf{C}\) 的一个 \(\mathrm{T}(\mathrm{X}) \otimes \mathbf{C}\)-自同构。存在 \(\mathrm{T}(\mathrm{X}) \otimes \mathbf{C}\) 作为两个复子丛直和的唯一分解：

\[
\mathbf {T} (\mathrm{X}) \otimes \mathbf {C} = \mathbf {T} ^ {\prime} (\mathrm{X}) \oplus \mathbf {T} ^ {\prime \prime} (\mathrm{X})\tag{1}
\]

使得 \(j\) 在 \(i\) 上与乘以 \(\mathbf{T}'(\mathbf{X})\) 一致，在 \(-i\) 上与乘以 \(\mathbf{T}''(\mathbf{X})\) 一致。\(p'\) 到 \(\mathbf{T}(\mathbf{X}) \otimes \mathbf{C}\) 的投影算子 \(\mathbf{T}'(\mathbf{X})\) 等于 \(\frac{1}{2}(1 - ij)\)；\(p'' = 1 - p'\) 到 \(\mathbf{T}(\mathbf{X}) \otimes \mathbf{C}\) 的投影算子 \(\mathbf{T}''(\mathbf{X})\) 等于 \(\frac{1}{2}(1 + ij)\)。复合映射

\[
\pi^ {\prime} \colon \mathrm{T} (\mathrm{X}) \rightarrow \mathrm{T} (\mathrm{X}) \otimes \mathbf {C} \xrightarrow {p ^ {\prime}} \mathrm{T} ^ {\prime} (\mathrm{X})
\]

\[
\pi^ {\prime \prime}: \mathrm{T} (\mathrm{X}) \rightarrow \mathrm{T} (\mathrm{X}) \otimes \mathbf {C} \xrightarrow {p ^ {\prime \prime}} \mathrm{T} ^ {\prime \prime} (\mathrm{X})
\]

是 R-同构；第一个将 j 映为乘以 i，第二个将 j 映为乘以 -i。

8.8.4（“\((p,q)\) 型的形式”）。保持前一号的假设和记号。设 F 是以 X 为基底的复向量丛，p 和 q 是两个 \(\geqslant0\) 的整数，且 \(n=p+q\)。设 \(\omega\) 是 X 的开集 U 上取值于 F 的 n 次微分形式；按照（8.8.2）将 \(\omega\) 与 \(\operatorname{Alt}_{\mathbf{C}}^{n}(\operatorname{T}(\mathbf{X})\otimes\mathbf{C};\mathbf{F})\) 的一个截面等同。若对所有 \(\omega\) 都有下式，则称 \((p,q)\) 为 \(x\in U\) 型：

\[
\omega_ {x} (v _ {1}, \dots , v _ {n}) = 0
\]

只要有 \(p + 1\) 个向量 \(v_i\) 属于 \(T_x'(X)\)，或有 \(q + 1\) 个向量 \(v_i\) 属于 \(T_x''(X)\)，该式就成立。若将 \(\bigwedge^n T_x(X) \otimes C\) 与各空间 \(\bigwedge^{m'} T_x'(X) \otimes \bigwedge^{m''} T_x''(X)\)（其中 \(m' + m'' = n\)）的直和等同（A, III, p. 85），并将 \(\omega_x\) 与一个 \(C\)-线性形式等同，该形式定义在 \(\bigwedge^n T_x(X)\) 上，则这等价于说，\(\omega_x\) 在 \(\bigwedge^{m'} T_x'(X) \otimes \bigwedge^{m''} T_x''(X)\) 上对 \((m', m'') \neq (p, q)\) 为零。

对于每一对满足 \((p,q)\) 的 \(\geqslant0\) 整数 \(p+q=n\)，存在且仅存在一个 \(\operatorname{Alt}^{p,q}(\mathrm{T}(\mathrm{X});\mathrm{F})\) 的复向量子丛 \(\operatorname{Alt}^{n}(\mathrm{T}(\mathrm{X});\mathrm{F})\)，使得其在 X 的开集 U 上的截面正是 X 上取值于 F 的 \((p,q)\) 型微分形式。向量丛 \(\operatorname{Alt}^{n}(\mathrm{T}(\mathrm{X});\mathrm{F})\) 是对于 \(\operatorname{Alt}^{p,q}(\mathrm{T}(\mathrm{X});\mathrm{F})\)、\(p\geqslant0\) 且 \(q\geqslant0\) 的向量丛 \(p+q=n\) 的直和。每个取值于 F 的 n 次微分形式 \(\omega\) 都唯一分解为

\[
\omega = \sum_ {p + q = n} \omega_ {p, q}
\]

其中 \(\omega_{p,q}\) 是一个 \((p,q)\) 型形式，称为 \((p,q)\) 的 \(\omega\) 型分量。若 \(\omega\) 属于 \(\mathbf{C}^h\) 类，其中 \(0 \leqslant h \leqslant k\)，则其各分量也具有同样性质。

例子

\begin{enumerate}
\def\labelenumi{\alph{enumi})}
\item
  n 次微分形式是 \((n, 0)\) 型，当且仅当它是 \(\mathbf{C}\)-多线性的。
\item
  设 \(\mathbf{F}'\)、\(\mathbf{F}''\) 和 \(\mathbf{F}\) 是三个复向量丛，且 \(\mathbf{F}' \times_{\mathbf{X}} \mathbf{F}'' \to \mathbf{F}\) 是一个 \(\mathbf{C}\)-双线性配对。若 \(\omega'\)（分别为 \(\omega''\)）是取值于 \((p',q')\)（分别为 \((p'',q'')\)）的 \(\mathbf{F}'\) 型（分别为 \(\mathbf{F}''\) 型）微分形式，则微分形式 \(\omega' \wedge \omega''\) 是 \((p' + p'',q' + q'')\) 型。
\item
  假设 F 是一个（实）向量丛的复化（8.8.2）。若 \(\omega\) 是取值于 F 的 \((p, q)\) 型微分形式，则其共轭 \(\overline{\omega}\)（8.8.2）是 \((q, p)\) 型。
\end{enumerate}

8.8.5. 除前述假设外，再假设 \(k \geqslant 1\)。以下三个条件等价：

\begin{enumerate}
\def\labelenumi{(\roman{enumi})}
\item
  对 X 的每个开集 U，以及 U 上每一对 \((\xi, \eta)\) 类复向量场 \(C^{k}\)，若对所有 \(\xi(x)\) 都有 \(\eta(x)\) 和 \(\mathrm{T}_{x}^{\prime}(\mathrm{X})\) 属于 \(x \in U\)，则对所有 \([\xi, \eta](x) \in \mathrm{T}_{x}^{\prime}(\mathrm{X})\) 都有 \(x \in U\)。
\item
  对 X 的每个开集 U，以及 U 上每一对 \((\xi, \eta)\) 类（实）向量场 \(C^{k}\)，有
\end{enumerate}

\[
[ \xi , \eta ] + j [ j \xi , \eta ] + j [ \xi , j \eta ] - [ j \xi , j \eta ] = 0.
\]

\begin{enumerate}
\def\labelenumi{(\roman{enumi})}
\setcounter{enumi}{2}
\tightlist
\item
  对 X 的每个开集 U，以及 U 上每个属于 \(\omega\) 类的 \((1,0)\) 型复微分形式 \(C^{k}\)，\((0,2)\) 的 \(d\omega\) 型分量为零。
\end{enumerate}

假设这些条件成立。设 \(\omega\) 是 X 的开集 U 上取值于复 Banach 空间 F、属于 \((p,q)\) 类且满足 \(\mathbf{C}^h\) 的 \(1\leqslant h\leqslant k\) 型微分形式。将 \(d'\omega\) 的 \(d''\omega\) 型（分别为 \((p+1,q)\) 型）分量记为 \((p,q+1))\)（分别为 \(d\omega\)）；此时 \(d\omega\) 的其他分量均为零，并且有

\[
d \omega = d ^ {\prime} \omega + d ^ {\prime \prime} \omega .\tag{2}
\]

若 \(k\geqslant 2\)，则有

\[
d ^ {\prime} (d ^ {\prime} \omega) = 0, d ^ {\prime \prime} (d ^ {\prime \prime} \omega) = 0, d ^ {\prime} (d ^ {\prime \prime} \omega) + d ^ {\prime \prime} (d ^ {\prime} \omega) = 0.\tag{3}
\]

采用 8.8.4 的记号，例 \(b\)，有

\begin{enumerate}
\def\labelenumi{(\arabic{enumi})}
\setcounter{enumi}{3}
\tightlist
\item
\end{enumerate}

\[
d ^ {\prime} \left(\omega^ {\prime} \wedge \omega^ {\prime \prime}\right) = d ^ {\prime} \left(\omega^ {\prime}\right) \wedge \omega^ {\prime \prime} + (- 1) ^ {p ^ {\prime} + q ^ {\prime}} \omega^ {\prime} \wedge d ^ {\prime} \left(\omega^ {\prime \prime}\right)\tag{4}
\]

\[
d ^ {\prime \prime} (\omega^ {\prime} \wedge \omega^ {\prime \prime}) = d ^ {\prime \prime} (\omega^ {\prime}) \wedge \omega^ {\prime \prime} + (- 1) ^ {p ^ {\prime} + q ^ {\prime}} \omega^ {\prime} \wedge d ^ {\prime \prime} (\omega^ {\prime \prime})\tag{5}
\]

结合 8.8.4 的例 c)，有

\[
d ^ {\prime} (\overline {{\omega}}) = \overline {{d ^ {\prime \prime} (\omega)}}.\tag{6}
\]

在公式 (4) 和 (5)（分别为 (6)）中，微分形式 \(\omega'\) 和 \(\omega''\)（分别为 \(\omega\)）均假定属于 \(\mathbf{C}^1\) 类。

8.8.6（“复流形的几乎复结构”）。设 \(X^c\) 是一个复解析流形，其底层实解析流形为 X（5.14.2.a)）。将丛 \(\mathrm{T}(\mathrm{X})\) 与复向量丛 \(\mathrm{T}(X^c)\) 的（实）底层向量丛等同，由此在 X 上赋予一个 \(C^\omega\) 类几乎复结构，称为由 X 上给定的复解析流形结构定义的几乎复结构。该几乎复结构满足 8.8.5 的条件 (i) 至 (iii)；特别地，8.8.5 的算子 \(d'\) 和 \(d''\) 有定义。若 \(f: X^c \to Y^c\) 是复解析流形的态射，则 \(f^*\) 与算子 \(d'\) 和 \(d''\) 可交换。

8.8.7. 保持 8.8.6 的假设和记号。令 \(\overline{\mathbf{X}}^c\) 为 \(\mathbf{X}^c\) 的共轭（5.14.2.b)）。通过对角映射 \(x \mapsto (x, x)\) 将 X 嵌入复流形 \(\mathbf{X}^c \times \overline{\mathbf{X}}^c\) 中，并令 \(g\) 为 \((x, y) \mapsto (y, x)\) 的反自同构 \(\mathbf{X}^c \times \overline{\mathbf{X}}^c\)。对 \((\mathbf{X}^c \times \overline{\mathbf{X}}^c, g)\) 是 X 的一个复化（5.14.8）。特别地，对每个 \(x \in \mathbf{X}\)，切空间 \(\mathrm{T}_x(\mathrm{X}^c \times \overline{\mathrm{X}}^c)\) 与 \(\mathrm{T}_x(\mathrm{X}) \otimes \mathbf{C}\) 等同；在此等同下，\(\mathrm{T}_x(\mathrm{X}^c \times \{x\})\) 对应于 \(\mathrm{T}_x'(\mathrm{X})\)，而 \(\mathrm{T}_x(\{x\} \times \overline{\mathrm{X}}^c)\) 对应于 \(\mathrm{T}_x''(\mathrm{X})\)。

8.8.8. 假设 \(r = \omega\)；X 上满足 8.8.5 的等价条件 (i) 至 (iii) 的任一 \(C^\omega\) 类几乎复结构，都是由 \(X\) 上唯一一个复解析流形结构定义的几乎复结构。\(^{(1)}\)

8.8.9（“全纯形式”）。恢复 8.8.6 的假设和记号，并设 F 是复 Banach 空间。\(\mathbf{X}^c\) 上取值于 F 的解析微分形式（换言之，复解析微分形式）仍称为全纯微分形式；将 \(\mathrm{T}(\mathbf{X}^c)\) 与 \(\mathrm{T}(\mathbf{X})\) 等同，就把它们与 X 上的（实解析）微分形式等同。对于 X 上取值于 F、属于 \(\omega\) 类且次数为 \(p\) 的微分形式 \(\mathbf{C}^k\)，其中 \(k \geqslant 1\)，它是全纯的充要条件是它属于 \((p, 0)\) 型且满足 \(d''\omega = 0\)。

设 \(\xi\) 是 X 上的向量场。若 \(\xi\) 是从 \(X^c\) 到 \(T(X^c) = T(X)\) 的复解析映射，换言之，是 \(X^c\) 上的复解析向量场，则称其为全纯的。在此情形，若 \(\omega\) 是 X 上取值于 F 的全纯微分形式，则形式 \(\theta_{\xi}.\omega\)（分别为 \(i(\xi)\omega\)）无论将 \(\xi\) 看作 \(X^c\) 上的复解析向量场、将 \(\omega\) 看作 \(\Omega^p(X^c; F)\) 的元素，还是将 \(\xi\) 看作 X 上的向量场、将 \(\omega\) 看作 \(\Omega^p(X; F)\) 的元素，都相同。此外，令 \(\xi' = \pi'(\xi)\) 和 \(\xi'' = \pi''(\xi)\) 分别为 \(\xi\) 在 \(T'(X)\) 和 \(T''(X)\) 中的分量（8.8.3）；有

\begin{enumerate}
\def\labelenumi{(\arabic{enumi})}
\setcounter{enumi}{6}
\tightlist
\item
\end{enumerate}
\[
\theta_ {\xi}. \omega = \theta_ {\xi^ {\prime}}. \omega , \quad \theta_ {\xi^ {\prime \prime}}. \omega = 0,\tag{7}
\]

\[
i (\xi) \omega = i (\xi^ {\prime}) \omega , i (\xi^ {\prime \prime}) \omega = 0.\tag{8}
\]

8.8.10. 例。--- 假设 \(X^{c}\) 是 \(C^{n}\) 的开子集，并记 \(z^{1},\ldots,z^{n}\) 为 \(X^{c}\) 上的坐标函数。对于 \(1\leqslant k\leqslant n\)，置 \(z^{k}=x^{k}+iy^{k}\)，其中 \(x^{k}\) 和 \(y^{k}\) 取实值。考虑由坐标系 \((\partial/\partial z^{k})\) 定义的 \(\mathrm{T}(\mathrm{X}^{c})\) 的切标架 \((z^{k})\)，以及由坐标系 \((\partial/\partial x^{k},\partial/\partial y^{k})\) 定义的 \(\mathrm{T}(\mathrm{X})\) 的切标架 \((x^{k},y^{k})\)（8.1.5）。有 \(j(\partial/\partial x^{k})=\partial/\partial y^{k}\)，且 \(\partial/\partial z^{k}\) 在映射 \(\pi^{\prime}\colon\mathrm{T}(\mathrm{X})=\mathrm{T}(\mathrm{X}^{c})\to\mathrm{T}^{\prime}(\mathrm{X})\)（8.8.3）下的像等于

\[
\frac {1}{2} (\partial / \partial x ^ {k} - i \partial / \partial y ^ {k});
\]

通常仍将其记为 \(\partial/\partial z^{k}\)。\(^{(2)}\) 该向量场 \(\partial/\partial z^{k}\) 的复共轭向量场记为 \(\partial/\partial\bar{z}^{k}\)；有

\[
\partial / \partial \bar {z} ^ {k} = \frac {1}{2} (\partial / \partial x ^ {k} + i \partial / \partial y ^ {k}).\tag{9}
\]

若 f 是 X 上取值于复 Banach 空间 F 的 \(C^{1}\) 类函数，则微分形式 \(d'f\) 和 \(d''f(8.8.5)\) 由下式给出

\[
d ^ {\prime} f = \sum_ {1 \leqslant k \leqslant n} \frac {\partial f}{\partial z ^ {k}} d z ^ {k}, d ^ {\prime \prime} f = \sum_ {1 \leqslant k \leqslant n} \frac {\partial f}{\partial \bar {z} ^ {k}} d \bar {z} ^ {k}.\tag{10}
\]

\(^{(1)}\) 若 X 是有限维的，则对于属于 \(C^k\) 类且满足 \(k = \dim X\) 的几乎复结构，上述结论仍然成立（参见 A. Newlander 和 L. Nirenberg, "Complex analytic coordinates in almost complex manifolds," Ann. of Math. LXX (1957), pp. 391--404）。

\(^{(2)}\) 须注意：将 \(\mathrm{T}(X^c)\) 与 \(\mathrm{T}(X)\) 典范等同，并不会把 \(\partial/\partial z^k\) 上的向量场 \(X^c\) 变为 X 上的向量场 \(\partial/\partial z^k\)，而是变为向量场 \(\partial/\partial z^k+\partial/\partial\bar z^k\)。然而，记作 \(\partial/\partial z^k\) 的两个向量场在全纯微分形式（8.8.9）上的取值相同；事实上，只有对这些形式才能应用 \(X^c\) 的向量场。

回顾有

\[
d f = d ^ {\prime} f + d ^ {\prime \prime} f
\]

\[
d z ^ {k} = d x ^ {k} + i d y ^ {k} \quad d \bar {z} ^ {k} = d x ^ {k} - i d y ^ {k}.
\]

若 \(\mathbf{A}=(i_{1},\ldots,i_{p})\) 是属于区间 \([1,n]\) 的严格递增整数序列，置

\[
d z ^ {A} = d z ^ {i _ {1}} \wedge \dots \wedge d z ^ {i _ {p}}
\]

并将 \(d\bar{z}^{\mathrm{A}}\) 的共轭记为 \(dz^{\mathrm{A}}\)。每个取值于 F 的 \(\omega\) 型微分形式 \((p,q)\) 都唯一地写成

\[
\omega = \sum_ {\mathrm{A}, \mathrm{B}} f _ {\mathrm{A}, \mathrm{B}} d z ^ {\mathrm{A}} \wedge d \bar {z} ^ {\mathrm{B}}\tag{11}
\]

其中 A（分别为 B）遍历由 \(p\) 中的 \(q\)（分别为 \([1, n]\)）个元素组成的严格递增序列集，\(f_{A,B}\) 是取值于 F 的函数。若 \(\omega\) 属于 \(\mathbf{C}^k\) 类，其中 \(k \geqslant 1\)，则函数 \(f_{A,B}\) 也具有同样的类，并且有

\begin{enumerate}
\def\labelenumi{(\arabic{enumi})}
\setcounter{enumi}{11}
\tightlist
\item
\end{enumerate}
\[
d ^ {\prime} \omega = \sum_ {A, B} d ^ {\prime} f _ {A, B} \wedge d z ^ {A} \wedge d \bar {z} ^ {B}\tag{12}
\]

\[
d ^ {\prime \prime} \omega = \sum_ {\mathbf {A}, \mathbf {B}} d ^ {\prime \prime} f _ {\mathbf {A}, \mathbf {B}} \wedge d z ^ {\mathbf {A}} \wedge d \bar {z} ^ {\mathbf {B}}.\tag{13}
\]

